% ============ §1 问题重述 ============
\section{问题重述}

\subsection{问题背景}

干燥是决定中药材成品品质的关键工序之一。含水率过高会导致药材霉变、有效成分降解；干燥速率过快则引起表面开裂、组织结构破坏。热风烘干因其设备简单、成本可控，是应用最广泛的干燥方式，其工艺通常分为两个阶段：预热平衡阶段（快速提升物料温度并保持较高环境湿度，防止表面过快失水）与恒温干燥阶段（维持稳定的温湿环境完成主体脱水）。然而，中药材干燥过程受物料内部传热传质规律的支配，扩散系数随含水率的剧烈变化、温度与含水率的双向耦合以及失水收缩等因素相互交织，使得工艺参数（烘房温度、湿度、时长）的选取高度依赖经验；传统试验优化模式成本高、周期长，亟需借助数理分析与数值仿真揭示干燥规律。

现有研究多采用有限差分/有限元对特定工况做个案仿真\cite{erzhi,moslae}，或采用薄层干燥动力学经验模型拟合干燥曲线\cite{food}。前者在竞赛级建模中缺乏解析洞察与验证深度，后者无法给出物料内部的时空分布。本文以谱配点方法为核心，构造兼具解析基准与数值精度的统一求解框架。

\subsection{问题描述}

某中药材形状大致呈圆柱形，长 25 cm、半径 2 cm；烘干开始时药材温度 28\,°C、水分浓度（干基含水率）2.55 kg/kg。烘房温度与水分浓度随时间的变化记录于附件 1（0--4 h，60 s 间隔）；烘干过程中药材半径的收缩数据记录于附件 2（0--72 h，30 min 间隔）。要求建立数学模型解决以下四个递进问题：

\begin{itemize}
\item \textbf{问题一}：建立预热平衡阶段（常物性，参数见附录 2）药材温度与水分浓度变化规律的数学模型，按表 1/表 2 格式给出 30 min 内 7 个时刻、5 个径向位置的结果，并将 1800 s 内每隔 1 s、径向每隔 0.1 cm 的完整结果保存至 result1.xlsx；
\item \textbf{问题二}：建立整个烘干过程（预热平衡 + 恒温干燥，变物性经验公式见附录 3）的数学模型，按表 3/表 4 格式给出 3 h 内每隔 0.5 h 的结果，并将完整结果保存至 result2.xlsx；
\item \textbf{问题三}：按烘干要求（药材各处水分浓度低于 0.15 kg/kg）确定烘干时长，按表 5 格式给出每隔 6 h 的水分浓度分布，并将完整结果保存至 result3.xlsx；
\item \textbf{问题四}：考虑实际烘干中因水分流失引起的尺寸变化（物性公式见附录 4，半径数据见附件 2），重新确定烘干时长，按表 6 格式给出结果并保存至 result4.xlsx。
\end{itemize}

所有结果保留四位小数。

% ============ §2 问题分析 ============
\section{问题分析}

\subsection{总体分析}

从传热传质与数值分析的交叉视角审视，本问题本质上是一个\textbf{圆柱轴对称域上的变物性耦合传热传质问题}，其难度由四个相互叠加的因素决定：

\begin{enumerate}
\item \textbf{强非线性}：附录 3 的扩散系数 $D=2.4\times10^{-3}e^{-0.45/C}e^{-3850/T_K}$ 随含水率 $C$ 从 2.55 降至 0.05 时跌落约 4 个数量级——这是干燥进入减速段、总时长达到 2--3 天的物理根源，也是数值求解必须采用隐式/自适应方法的直接原因；
\item \textbf{双向耦合}：含水率改变密度、比热容、导热系数（热惯性），温度通过 Arrhenius 项改变扩散系数（迁移速率）；
\item \textbf{两阶段时变边界}：烘房环境在 0--4 h 内由数据给定并逐渐进入平台期，恒温段边界条件需要数据驱动的外推；
\item \textbf{移动边界}：问题四中半径随时间收缩（附件 2），求解域随时间变化。
\end{enumerate}

先做无量纲量分析以预判物理行为（这也是数值参数选取的依据）：热扩散率 $\alpha=k/(\rho c_p)\approx1.69\times10^{-7}$ m$^2$/s，Biot 数 $\mathrm{Bi}_T=hR/k=1.39$，Fourier 数 $\mathrm{Fo}_T(1800\,\mathrm{s})=0.76$——温度场在 30 min 内接近 4.4 个热时间常数；水分扩散 Fourier 数 $\mathrm{Fo}_m(1800\,\mathrm{s})\approx0.022$——中心含水率半小时内几乎不变，预热阶段“锁水”的工艺意图由此得到定量刻画。局部 Luikov 数 $\mathrm{Lu}=\alpha/D$ 由湿态的约 10 升至干态的约 $10^5$，表明热湿耦合强度随干燥进程急剧衰减（此观察将用于解释层与灵敏度分析）。

\subsection{问题一分析}

类型：常物性轴对称瞬态传热传质（$D$ 仅随 $C$ 变化，弱耦合）。关键难点有二：其一，4 位小数的表格精度要求数值格式与求解器经得起验证；其二，环境数据为离散采样，边界条件需要光滑且保真的处理。策略：将烘房环境辨识为一阶惯性环节（物理上对应加热系统惯性），对指数边界导出温度场的 Duhamel 闭式解析解（精确、可作基准）；水分场用 Chebyshev 谱配点求解，并构造“解析解互证 + Bessel 级数基准 + 制造解验证”的验证链。

\subsection{问题二分析}

类型：变物性双向耦合瞬态传热传质。新增物理条件：物性随 $C$、$T$ 变化（附录 3）；环境条件需覆盖 2--3 天（附件 1 只给 4 h，需外推恒温段）。策略：沿用问题一的谱配点框架，将物性函数代入；边界条件两阶段拼接（0--4 h 辨识模型，之后平台外推）；温度--水分耦合在边界消元处联立求解。验证重点：干燥锋面的谱分辨率（此为本问最易出错的数值环节）。

\subsection{问题三分析}

类型：问题二模型的判据扫描（决策条件）。判据 $\max C(r,t)<0.15$：由水分沿半径单调递减（扩散外流），最严点在中心。直接复用问题二的解，在 60 s 分辨率上扫描判据跨越时刻。关键讨论点：跨越处中心曲线极平缓，$t_{end}$ 对数值精度敏感——需给出收敛序列与交叉验证区间。

\subsection{问题四分析}

类型：移动边界（收缩）变物性耦合问题。新增物理条件：$R=R(t)$（附件 2 实测）；物性换附录 4；求解域随时间收缩。策略：任意拉格朗日--欧拉（ALE）表述——网格随体运动，干基含水率的物质守恒给出极简方程（对流、膨胀、密度变化项精确抵消），几何守恒律（GCL）保证离散一致性；与问题三对比量化收缩效应。
